% corpus_supplement: Erdős problem #1082
% source: https://www.erdosproblems.com/latex/1082
% fetched_at: 2026-07-14T18:48:27Z
% payload_sha256: 3cfcb1ae0a80151ba6c1ab2d613f173c43e5b460e988bd0ab7097b505b09f844
% statement/remarks/references extracted by erdos_ingest.extract_source_page;
% rendered by erdos_ingest.render_tex — do not edit
\documentclass{article}
\usepackage{amsmath, amssymb, amsthm}
\usepackage{hyperref}
\usepackage{geometry}
\geometry{margin=1in}

\title{Erd\H{o}s Problem \#1082}
\date{Last updated: 2025-10-17}
\author{Source: erdosproblems.com}

\begin{document}
\maketitle

\noindent\textbf{Status:} OPEN \quad \textbf{No prize}

\noindent\textbf{Tags:} geometry, distances

\medskip
\noindent\textbf{Problem Statement:}

Let $A\subset \mathbb{R}^2$ be a set of $n$ points with no three on a line. Does $A$ determine at least $\lfloor n/2\rfloor$ distinct distances? In fact, must there exist a single point from which there are at least $\lfloor n/2\rfloor$ distinct distances?

\bigskip
\noindent\textbf{Remarks:}

A conjecture of Szemer\'{e}di, who proved this with $n/2$ replaced by $n/3$. More generally, Szemer\'{e}di gave a simple proof that if there are no $k$ points on a line then some point determines $\gg n/k$ distinct distances (a weak inverse result to the distinct distance problem [89]).

This is a stronger form of [93]. The second question is a stronger form of [982].

Szemer\'{e}di's proof is unpublished, but given in \cite{Er75f}.

In \cite{Er75f} Erd\H{o}s asks whether, given $n$ points in $\mathbb{R}^3$ with no three on a line, do they determine $\gg n$ distances? Altman proved the answer is yes if the points form the vertices of a convex polyhedron (see [660] for a stronger form of this), and Szemer\'{e}di proved the answer is yes if there are no four points on a plane.

The second stronger question has a negative answer: there is {IMAGE=1082,a construction} of $8$ points such that each point has exactly three distinct distances to the others. This first appeared in the literature in a paper of Erd\H{o}s and Fishburn \cite{ErFi97b}, where they credit the construction to Harborth. This configuration is studied in detail by Fishburn \cite{Fi02}.

(This construction was later found by DeepMind. An earlier construction was also given by Xichuan in the comments, who gave a set of $42$ points in $\mathbb{R}^2$, with no three on a line, such that each point determines only $20$ distinct distances.)

\bigskip
\noindent\textbf{References:}

[Er75f] Erd\H{o}s, Paul, On some problems of elementary and combinatorial geometry. Ann. Mat. Pura Appl. (4) (1975), 99-108.
 
 [ErFi97b] Erd\H{o}s, Paul and Fishburn, Peter, Distinct distances in finite planar sets. Discrete Math. (1997), 97--132.
 
 [Fi02] Fishburn, Peter C., A remarkable eight-point planar configuration. Discrete Math. (2002), 103--122.

\bigskip
\noindent\small{Source: \url{https://www.erdosproblems.com/1082}}
\end{document}
